En ÍndicePartidas, una unidad es una cuenta que paga una suscripción.
El objetivo de este cálculo es estimar cuánto dinero aporta esa cuenta
y cuánto tarda en recuperar lo que costó conseguirla.

\paragraph{Cuánto deja una cuenta al mes}
Los precios propuestos en \texttt{Pricing.md} son \$9.990 para el plan
Independiente y \$24.990 para el plan Pyme, ambos antes de IVA \hip.
Si el 70\,\% de las cuentas usa el primer plan y el 30\,\% usa el segundo,
el ingreso mensual promedio por cuenta (ARPA) es de \$14.490 \hip.
Este reparto 70/30 es un supuesto deliberadamente conservador para evaluar
la viabilidad mínima con predominio del plan de menor barrera económica;
no obstante, las entrevistas de campo (\texttt{E-02}, \texttt{E-03}, \texttt{E-06})
muestran mayor dolor operativo y disposición al pago en el segmento Pyme,
por lo que si el piloto concentra más cuentas en ese plan, el ARPA aumentará
favorablemente \hip.
Estimamos costos directos de \$2.500 por cuenta Independiente y \$6.000 por
cuenta Pyme para operar el servicio, dar soporte y procesar el cobro.
Este valor desglosa comisiones de pasarela de pago ($\sim$3,2\,\% por cobro recurrente),
costo de infraestructura cloud prorrateado entre cuentas activas (\$1.150 a \$2.000
por cuenta) y soporte técnico promedio prorrateado de 6 a 18 minutos mensuales por
cuenta activa (0,1 a 0,3 horas valorizadas a \$10.000 la hora de dedicación del equipo) \hip.
Con esa misma mezcla, atender una cuenta cuesta en promedio \$3.550 al mes
\hip. La diferencia es una \textbf{contribución mensual de \$10.940 por
cuenta}, equivalente al 75,5\,\% del ingreso \hip.

\paragraph{Cuánto cuesta conseguir una cuenta}
Si gastamos \$120.000 entre captación y atención de la prueba gratuita,
y cuatro cuentas empiezan a pagar, cada cuenta cuesta \$30.000 de adquirir
(CAC) \hip. Esta cifra estima entre tres y cinco contactos gestionados en el período
de prueba, con dos horas de coordinación por contacto a \$10.000 la hora de dedicación
(\$60.000 a \$100.000), más tiempo de soporte durante los 30 días (\$20.000 a \$40.000),
utilizando \$120.000 como valor central del rango \hip. Al dividir \$30.000 por
los \$10.940 que deja cada mes, obtenemos un \textbf{plazo de recuperación de 2,74 meses}
desde que la cuenta empieza a pagar \hip.

\paragraph{Cuánto aporta antes de cancelar}
Si cada mes cancela el 10\,\% de las cuentas que ya pagan
(\textit{churn}) y esa proporción se mantiene, una cuenta permanece activa
unos 10 meses en promedio \hip. En ese tiempo aporta cerca de \$109.400
(LTV), es decir, 3,65 veces lo que cuesta conseguirla (LTV/CAC) \hip.
Con esos supuestos, una cuenta cubre su costo de adquisición. Para saber
si el negocio completo es rentable todavía hay que considerar los costos
fijos de desarrollo y administración.

\paragraph{Qué pasa si menos cuentas empiezan a pagar}
Si con los mismos \$120.000 conseguimos solo dos cuentas pagadas, el costo
por cuenta sube a \$60.000, la relación LTV/CAC baja a 1,82 y la recuperación
tarda 5,48 meses \hip. En el piloto registraremos cuánto gastamos en
captación y soporte, cuántas cuentas pasan de la prueba al pago y cuántas
cancelan cada mes. Si convertir una cuenta a pago resulta más costoso de
lo esperado, probaremos una configuración inicial guiada y mediremos si
las conversiones adicionales compensan el tiempo de soporte.
